% answers_short.tex —— 2015 年·天津卷·文科 · 简短答案速查
% 由 tools/make_short_answers.py 自 answers.tex 的 \gksolhead/\gkanswer 提取，
% 并与 2015/2015_数学真题_天津文_答案要点.md 逐题比对后生成，勿手改（改源 MD 后重跑）。

\begin{gknotice}
\textbf{使用说明}：本节只给最终答案便于快速核对，完整推导见解析版。答案为编写者独立推导，\textbf{非官方标准答案}。
\end{gknotice}

\gkpart{选择题}
\begin{center}\begin{tabular}{ccccccccc}
\toprule
题号 & 1 & 2 & 3 & 4 & 5 & 6 & 7 & 8 \\
\midrule
答案 & B & C & C & A & D & A & B & A \\
\bottomrule
\end{tabular}\end{center}

\gkpart{填空题}
\begin{itemize}[leftmargin=2.6em]
  \item[\textbf{9.}] \(-\mathrm{i}\)
  \item[\textbf{10.}] \(\dfrac{8\pi}{3}\ \mathrm{m}^3\)
  \item[\textbf{11.}] \(3\)
  \item[\textbf{12.}] \(4\)
  \item[\textbf{13.}] \(\dfrac{29}{18}\)
  \item[\textbf{14.}] \(\dfrac{\sqrt{\pi}}{2}\)
\end{itemize}

\gkpart{解答题}
\begin{itemize}[leftmargin=2.6em]
  \item[\textbf{15.}] (1) 甲 3 人、乙 1 人、丙 2 人；(2) ① 15 个基本事件；② \(\dfrac{3}{5}\)
  \item[\textbf{16.}] (1) \(a=8\)，\(\sin C=\dfrac{\sqrt{15}}{8}\)；(2) \(\dfrac{\sqrt{15}-7\sqrt{3}}{16}\)
  \item[\textbf{17.}] (1) 成立（\(E\)，\(F\) 分别为 \(BC\)，\(A_1C\) 中点，\(EF\) 为 \(\triangle CBA_1\) 的中位线，\(EF\parallel A_1B\)，故 \(EF\parallel\) 平面 \(A_1B_1BA\)）；(2) 成立（\(AB=AC\) 且 \(E\) 为 \(BC\) 中点得 \(AE\perp BC\)，又 \(AA_1\perp\) 平面 \(ABC\) 得 \(AA_1\perp BC\)，故 \(BC\perp\) 平面 \(AEA_1\)，从而平面 \(AEA_1\perp\) 平面 \(BCB_1\)）；(3) \(30^{\circ}\)
  \item[\textbf{18.}] (1) \(a_n=2^{n-1}\)，\(b_n=2n-1\)；(2) \(S_n=(2n-3)\cdot2^{n}+3\)
  \item[\textbf{19.}] (1) \(2\)；(2) ① \(\lambda=\dfrac{7}{8}\)；② \(\dfrac{x^2}{5}+\dfrac{y^2}{4}=1\)
  \item[\textbf{20.}] (1) \(f(x)\) 在 \((-\infty,1)\) 上单调递增，在 \((1,+\infty)\) 上单调递减；(2) 成立（\(P(4^{\frac{1}{3}},0)\)，切线 \(g(x)=-12x+12\cdot4^{\frac{1}{3}}\)，\(g(x)-f(x)=x^4-16x+12\cdot4^{\frac{1}{3}}\geqslant0\) 恒成立，等号仅当 \(x=4^{\frac{1}{3}}\)）；(3) 成立（由切线不等式得 \(x_2<4^{\frac{1}{3}}-\dfrac{a}{12}\)，又 \(x_1>\dfrac{a}{4}\)，相减得 \(x_2-x_1<4^{\frac{1}{3}}-\dfrac{a}{3}\)，故所证不等式成立）
\end{itemize}
